\documentclass[11pt,a4paper]{article}
\usepackage[T1]{fontenc}
\usepackage[utf8]{inputenc}
\usepackage[ngerman]{babel}
\usepackage[margin=2.5cm]{geometry}
\usepackage[hidelinks]{hyperref}
\usepackage{array}
\setlength{\parindent}{0pt}
\setlength{\emergencystretch}{2em}
\renewcommand{\arraystretch}{1.25}
\begin{document}
\shorthandoff{"}
\begin{center}
  {\Large\textbf{Webtechnologien -- Übungsblatt 2}}
\end{center}
\section*{Aufgabe 4: Glass Price}
\textbf{Aufgabe:} In \texttt{calculate.js} den Preis einer Glasbestellung
berechnen und in einem neuen Preisfeld anzeigen.
\texttt{glassOrder.html} darf nicht geändert werden.

\medskip
\textbf{Ergebnis:}
\begin{itemize}
  \item Doppelte Überschrift entfernt:
    \texttt{document.querySelector("h1 + p").remove()}.
    \texttt{h1 + p} wählt den Absatz direkt nach dem \texttt{h1}.
  \item Preisfeld: neues \texttt{input} (\texttt{type="text"}, ID und Name
    \texttt{price}, \texttt{readOnly}) mit \texttt{label} \glqq price\grqq{}
    (\texttt{htmlFor = "price"}). Beide werden mit \texttt{append} ans Ende
    des Fieldsets \glqq Glass Specification\grqq{} gehängt, also nach der
    Legende \glqq Price\grqq{} an der Stelle des Kommentars.
  \item \texttt{calculate(width, height, thickness)}: Das Farbfeld liefert
    \texttt{\#rrggbb}. Rot, Grün und Blau werden mit \texttt{substring}
    herausgelöst. Sind alle drei gleich, ist das Glas ein Grauton
    ($b \cdot h \cdot d / 50$), sonst farbig ($b \cdot h \cdot d / 32$).
    Schwarz und Weiß zählen ebenfalls als Grauton.
  \item Einheiten: Breite und Höhe sind laut Formular in cm, die Dicke
    (4, 6, 8, 12) in mm. Die Dicke wird daher durch 10 geteilt.
  \item \texttt{update()}: prüft Breite, Höhe und die Dicken-Gruppe mit
    \texttt{checkValidity()}. Geprüft werden \texttt{required},
    \texttt{min="0"}, \texttt{max="200"} und \texttt{step="1"}.
    Ein Radio mit \texttt{required} ist gültig, sobald irgendein Radio
    derselben Gruppe gewählt ist; deshalb genügt
    \texttt{thicknessInputs[0]}. Gültig: Preis mit zwei Nachkommastellen,
    sonst \glqq -\grqq.
  \item Events: \texttt{input} auf Breite, Höhe und Farbe (feuert bei jeder
    Änderung, auch beim Tippen), \texttt{change} auf jedem Radio.
    \texttt{update()} läuft einmal beim Laden und zeigt \glqq -\grqq.
\end{itemize}

\textbf{Hinweis:} Für den Preis werden nur die preisrelevanten
Pflichtfelder geprüft (Breite, Höhe, Dicke). Die Kontaktdaten ändern den
Preis nicht. Das Farbfeld hat immer einen Wert (Standard \texttt{\#000000}).

\medskip
\textbf{Test:} In Chrome, \texttt{glassOrder.html} direkt geöffnet:
\begin{center}
\begin{tabular}{|>{\raggedright\arraybackslash}p{6.3cm}|>{\raggedright\arraybackslash}p{5.2cm}|>{\raggedright\arraybackslash}p{2.2cm}|}
\hline
\textbf{Eingabe} & \textbf{Rechnung} & \textbf{Preisfeld} \\ \hline
$100 \times 50$\,cm, 4\,mm, \texttt{\#000000} & $100 \cdot 50 \cdot 0{,}4 / 50$ & \texttt{40.00} \\ \hline
$100 \times 50$\,cm, 4\,mm, \texttt{\#ff0000} & $100 \cdot 50 \cdot 0{,}4 / 32$ & \texttt{62.50} \\ \hline
$100 \times 50$\,cm, 12\,mm, \texttt{\#808080} & $100 \cdot 50 \cdot 1{,}2 / 50$ & \texttt{120.00} \\ \hline
Breite 300 & über \texttt{max="200"} & \texttt{-} \\ \hline
Breite leer & \texttt{required} verletzt & \texttt{-} \\ \hline
nur Breite eingegeben & Höhe und Dicke fehlen & \texttt{-} \\ \hline
\end{tabular}
\end{center}
Die HTML-Datei ist unverändert (bytegleich mit der Vorlage geprüft).
Das Formular wurde nicht abgeschickt, da es über \texttt{mailto} sendet.

\medskip
{\small \textbf{Hilfsmittel:} Claude Code (Claude Opus 5.5) für Erstellung, Prüfung und Dokumentation.}
\end{document}
